\section{Lifting Real Robot Data to 4D and Pretraining}
\label{app:pretraining}

\subsection{Data, Partitions, and Time}
\label{app:droid}
The pretraining corpus is DROID \cite{droid}, recorded at 15\,Hz with two exterior
cameras and one wrist camera, together with the 3D point tracks and optimized camera
calibrations that PointWorld released for it \cite{pointworld}. We associate each
released clip with its raw DROID episode by requiring exact agreement of the joint
sequence, end-effector poses, gripper states, commands, and camera identity, and we
exclude clips whose capture clock is not monotonic or whose association is
ambiguous. The model sees the first exterior camera. The training unit is a causal
observation window with its aligned future and action labels. We assign episodes
to training, development, and held-out partitions (90/5/5, with a fixed random
seed) by their original identifier before any filtering, so overlapping clips of an
episode cannot cross partitions. Table~\ref{tab:pretraining_data} reports the yield
of every stage and counts episodes, windows, actions, and valid labels separately.

% Planned results: DROID lifting statistics and stage yield (App. E). PW0 manifest facts for reference (confirm before submission): 42,935 released episodes; 42,087 trainable; 598,889 main-camera windows (train 37,874 ep. / 538,740 win.; val 4,213 / 60,149); window action validity 96.88%.
\begin{table}[t]
\centering\footnotesize
\caption{\tablelang{\textbf{DROID lifted to 4D.} Yield of each stage. Windows are causal observation windows with aligned
labels. Exogenous windows show target motion before the first robot--target contact (Appendix~\ref{app:lifting}). The last
two rows compare all windows with valid windows.}{\textbf{提升为 4D 的 DROID。}各阶段的保留量。窗口指带有对齐标签的因果观测窗口。
外源窗口指在机器人与目标首次接触前目标已在运动的窗口（附录~\ref{app:lifting}）；最后两行比较全部窗口与有效窗口。}}
\label{tab:pretraining_data}
\setlength{\tabcolsep}{3.4pt}
\begin{tabular}{@{}lccc@{}}
\toprule
 & \tablelang{Train}{训练} & \tablelang{Dev.}{开发} & \tablelang{Held-out}{留出} \\
\midrule
\tablelang{Episodes with PointWorld tracks}{含 PointWorld 轨迹的 episode} & \resultplaceholder & \resultplaceholder & \resultplaceholder \\
\quad\tablelang{associated to raw DROID}{与原始 DROID 对齐} & \resultplaceholder & \resultplaceholder & \resultplaceholder \\
\quad\tablelang{one rigid target named and boxed}{指令只指定一个刚体且已检测} & \resultplaceholder & \resultplaceholder & \resultplaceholder \\
\tablelang{Windows}{窗口} & \resultplaceholder & \resultplaceholder & \resultplaceholder \\
\quad\tablelang{exogenous (\%)}{外源（\%）} & \resultplaceholder & \resultplaceholder & \resultplaceholder \\
\quad\tablelang{with causal pose input}{含因果位姿输入} & \resultplaceholder & \resultplaceholder & \resultplaceholder \\
\tablelang{Valid actions (\%)}{有效动作（\%）} & \resultplaceholder & \resultplaceholder & \resultplaceholder \\
\tablelang{Valid SE(3) labels (horizons, \%)}{有效 SE(3) 标签（时域，\%）} & \resultplaceholder & \resultplaceholder & \resultplaceholder \\
\tablelang{Valid 3D flow labels (points, \%)}{有效 3D 流标签（点，\%）} & \resultplaceholder & \resultplaceholder & \resultplaceholder \\
\tablelang{Speed, all / valid (cm/s)}{速度，全部 / 有效（cm/s）} & \resultplaceholder[0.0]/\resultplaceholder[0.0] & \resultplaceholder[0.0]/\resultplaceholder[0.0] & \resultplaceholder[0.0]/\resultplaceholder[0.0] \\
\tablelang{Occluded, all / valid (\%)}{遮挡，全部 / 有效（\%）} & \resultplaceholder[00]/\resultplaceholder[00] & \resultplaceholder[00]/\resultplaceholder[00] & \resultplaceholder[00]/\resultplaceholder[00] \\
\bottomrule
\end{tabular}
\end{table}


\paragraph{Time.} All times are capture timestamps, and we do not relabel frames to
a nominal rate. The history offsets lie on a 50\,ms grid that the 15\,Hz frames do
not hit, so we resample the causal inputs at the offsets exactly as at deployment
(Appendix~\ref{app:tokens}), between estimates that had been emitted by the query
time and without extrapolation. Target labels interpolate the offline trajectory in
the same way, linearly in position and along the geodesic in rotation, and not
across gaps longer than two source periods. They therefore lie on the horizons of
the main model. Flow labels keep PointWorld's own record times, \fillin{0.13} to
\fillin{1.13}\,s ahead, without interpolation. The world readouts embed each
horizon's offset $\delta_h$ in physical time (Appendix~\ref{app:world}), so the
scene queries are supervised at these offsets during pretraining and on the 50\,ms
grid afterwards.

\paragraph{Actions.} We express robot states and Cartesian commands in the robot
base frame after the calibrated flange-to-TCP transform. Command labels follow
Appendix~\ref{app:action}. Each label is the increment from the pose commanded for
the preceding application time to the pose commanded for the current one, with
DROID's commanded poses interpolated to the application times $\tau_t+\ell+j\Delta_a$
of the sampled latency. We admit a label if its values are finite and its reference
state was available causally. The robot's execution limits are not an admission
rule, because a per-axis limit of 5\,cm per command would discard \fillin{XX}\% of
the commands in the fastest quartile of motion and bias the base toward slow
demonstrations. We mask only the steps at which DROID's collector issued no command
(\fillin{3.1}\% of window slots). Action and flow validity are independent, so an
invalid action does not remove the flow labels of its window.

\subsection{Target and Its Pose}
\label{app:droid_target}
DROID provides no object poses, and its instructions name the manipulated object
without marking it. We take as the target the object named in the instruction,
i.e., the direct object of the main verb extracted by a dependency parser (``the
mug'' in ``put the mug in the sink''), with one persistent instance per episode.
Grounding DINO \cite{groundingdino} boxes this phrase in the first frame in which it
scores at least 0.35, and SAM~2.1 \cite{sam2} tracks its mask. For every frame with
a non-empty mask, we prompt WildDet3D \cite{wilddet3d} with the tight 2D box, the
RGB frame, and the fixed intrinsics. It predicts a metric 3D box from the single
image with its own geometry estimate and no external depth, and PointWorld's
optimized extrinsics map the box to the base frame. We accept a box when the box
centers from the two exterior cameras lie within 5\,cm of each other in the base
frame. We exclude episodes whose instruction names no object, several objects, or a
non-rigid one, as well as episodes without an accepted box. An empty mask counts as
a missing observation, and we do not fill it with the previous pose.

\paragraph{Body frame.} A 3D box defines its frame only up to the symmetries of the
box, and the detector may permute or flip its axes from one frame to the next. We
therefore map each new box to the equivalent frame (an axis permutation and sign
flips that preserve its extents) nearest to the previous valid estimate, which fixes
one body frame per episode. Rotation steps above $45^\circ$ between consecutive
frames that remain after this step, in \fillin{X.X}\% of frames, are marked invalid
instead of being smoothed. As in simulation, rotation about the symmetry axis of a
symmetric target is excluded from the loss (Sec.~\mainref{sec:predict}).

\subsection{Causal Inputs and Offline Labels}
\label{app:lifting}
\begin{figure*}[t]
\centering
\figplaceholder[1.25in]{\textwidth}{Fig.~S2 --- lifting a real robot episode to metric 4D (schematic)}{%
\textbf{Left (shared setup):} the instruction with its object phrase underlined (``put \underline{the mug} in the sink'') and the
Grounding DINO box on the first frame of both exterior views, whose WildDet3D box centers agree in the base frame; a small inset shows
PointWorld's optimized camera pose next to the raw one. \textbf{Top row (causal input path, \textcolor{obscol}{blue}):} frames arrive
left to right; SAM~2 propagates the mask forward and WildDet3D lifts each masked frame to a metric 3D box; emitted poses are triads
stamped with their availability time; an empty mask shows an explicit ``missing'' token. No arrow points backwards in time.
\textbf{Bottom row (offline label path, \textcolor{predcol}{orange}):} masks propagated forward and backward, per-frame boxes mapped to
one body frame, the two pose tracks fused by the robust smoother with an abrupt event preserved (the grasp), cumulative SE(3) labels
$(\vect d_h,\boldsymbol\phi_h)$ on the 50\,ms horizons; a second branch picks, in each grid patch, the PointWorld track nearest the
patch center and draws its 3D displacement at PointWorld's record times, with invalid points greyed out. A vertical dashed line marks
the first robot--target contact, after which the masked-label control drops world labels. \textbf{Right:} one training sample =
causal inputs (top) + masked labels (bottom). Use one real DROID frame sequence, vector overlays only.}
\caption{Lifting DROID to metric 4D supervision. The object named in the instruction
is boxed and tracked. The inputs reproduce deployment with forward-only tracking,
whereas the labels use the whole episode, with bidirectional tracking and smoothing,
and PointWorld's 3D tracks. Labels do not rewrite inputs, and invalid labels mask
their losses.}
\label{fig:lifting}
\end{figure*}

\paragraph{Causal inputs.} SAM~2.1 (Hiera-small) propagates the mask causally, and
WildDet3D lifts every masked frame in full precision, so that the estimate does not
depend on the batch composition. Each query processes only the newly available
frames and appends its estimate without rewriting earlier ones. Empty masks,
rejected boxes, and rotation jumps set the input validity to false, and training
uses these naturally occurring missing states without synthetic dropout. The real
robot uses the same path for objects without a CAD model or reference views
(Appendix~\ref{app:tracker}).

\paragraph{Offline labels.} The labels use the full episode and the same target
phrase. SAM~2 propagates the mask forward and backward, and we accept frames whose
two masks reach an IoU of at least 0.5. The boxes from both directions enter one
robust smoother, a Huber cost with pose residuals scaled by 5\,mm and $3^\circ$ and
velocity-change residuals weighted by 0.1. The velocity term is omitted across gaps
and across velocity jumps above 0.2\,m/s or 1\,rad/s, so that abrupt events are
preserved. Twenty manifold Gauss--Newton iterations optimize the cost, and frames
whose two boxes disagree by more than 20\,mm or $15^\circ$ are invalid. The target
labels are $\tilde{\vect p}_{t+h}-\tilde{\vect p}_t$ and
$\Log(\tilde R_{t+h}\tilde R_t^\top)^\vee$ within the offline trajectory. Because
they are not offline futures minus causal estimates, they do not absorb corrections
of the current pose. A horizon is valid only if its reference and future poses are
valid and share identity and body frame.

\paragraph{Scene flow.} PointWorld's tracks are already expressed in the base frame.
For each point of the grid, the label follows the released track whose initial
projection lies nearest to the point within its patch, provided that the track is
initially visible, has a valid positive depth inside the image, and agrees with the
depth at its nearest pixel within 20\,mm. Patches without such a track are invalid,
and we do not fill them from neighbors. The label is the track's displacement from
its initial position at each record time, and its visibility is PointWorld's. An
occluded, out-of-image, or depth-invalid future point masks its displacement loss
and supervises only its visibility. The selected tracks lie \fillin{0.5}\,px from
the patch center at the median (\fillin{9}\,px at the 95th percentile).

\paragraph{Contact and exogenous motion.} The first robot--target contact is the
first frame at which the target box comes within 1\,cm of the gripper's finger
geometry from the robot's kinematics. A window is \emph{exogenous} if the target
moves by more than 1\,cm or $5^\circ$ within its 400\,ms horizon before that contact,
and Table~\ref{tab:pretraining_data} reports their share. The masked-label control
of Table~\ref{tab:supp_pretraining} removes all world labels after the first contact
and keeps the action loss. It retains \fillin{XX}\% of the low-data gain of
pretraining, so \fillin{most of this gain does not come from predicting how the
robot moves the target}.

\paragraph{Validity bias.} PointWorld's tracks cover static surfaces more completely
than fast-moving ones, and accepted boxes favor unoccluded, well-textured targets.
Table~\ref{tab:pretraining_data} compares the target speed and the occlusion between
all windows and valid windows. The label-source control of
Table~\ref{tab:supp_pretraining} replaces the tracks by flow lifted from 2D
correspondences \cite{dpflow} with monocular metric depth \cite{moge2}, which covers
more of the image.

\subsection{Label Quality}
\label{app:label_quality}
% Planned results: controlled DROID-base comparisons and data efficiency (App. E). Simulation post-training from each initialization; DROID held-out forecast error of each base. Real-robot effects of the base: Tables 5, S-real.
\begin{table*}[t]
\centering\small
\caption{\tablelang{\textbf{Controlled comparison of initializations.} Each row post-trains the simulation model from its
own initialization on 10/25/50/100\% of the \bench demonstrations (held-out SR, \%, three post-training seeds). The
generic initialization is the one \ours uses. The last columns give each base's forecast error on held-out DROID episodes
at the last horizon, for target translation (mm) and rotation ($^\circ$) against the smoothed labels and for the flow
endpoint (mm) against PointWorld tracks. Control bases are pretrained at pilot scale against a world--action base at the same
budget (Appendix~\ref{app:scale}).}{\textbf{初始化的受控比较。}每行从各自的初始化出发，用 10/25/50/100\% 的 \bench 示范后训练仿真模型
（留出 SR，\%，三个后训练种子）；通用初始化即 \ours 所用的初始化。最后几列为各基座在留出 DROID episode 上最后一个时域的预测误差：
目标平移（mm）与旋转（$^\circ$，相对平滑标签），以及相对 PointWorld 轨迹的流端点误差（mm）。对照基座按试点规模预训练，并与相同预算的世界—动作基座比较（附录~\ref{app:scale}）。}}
\label{tab:supp_pretraining}
\setlength{\tabcolsep}{4.6pt}
\begin{tabular}{@{}lccccccc@{}}
\toprule
 & \multicolumn{4}{c}{\tablelang{\bench SR, fraction of demonstrations}{\bench SR，示范比例}} & \multicolumn{3}{c}{\tablelang{DROID held-out}{DROID 留出}} \\
\cmidrule(lr){2-5}\cmidrule(lr){6-8}
\tablelang{Initialization}{初始化} & 10\% & 25\% & 50\% & 100\% & mm & $^\circ$ & \tablelang{flow (mm)}{流（mm）} \\
\midrule
\tablelang{Generic visual only (\ours)}{仅通用视觉初始化（\ours）} & \resultplaceholder & \resultplaceholder & \resultplaceholder & \resultplaceholder & -- & -- & -- \\
\tablelang{Action-only base, same episodes}{仅动作基座，相同 episode} & \resultplaceholder & \resultplaceholder & \resultplaceholder & \resultplaceholder & -- & -- & -- \\
\tablelang{World--action base, flow from monocular depth}{世界—动作基座，单目深度提升的流} & \resultplaceholder & \resultplaceholder & \resultplaceholder & \resultplaceholder & \resultplaceholder & \resultplaceholder & \resultplaceholder \\
\tablelang{World--action base, world labels masked after contact}{世界—动作基座，接触后屏蔽世界标签} & \resultplaceholder & \resultplaceholder & \resultplaceholder & \resultplaceholder & \resultplaceholder & \resultplaceholder & \resultplaceholder \\
\tablelang{World--action base, pilot scale}{世界—动作基座，试点规模} & \resultplaceholder & \resultplaceholder & \resultplaceholder & \resultplaceholder & \resultplaceholder & \resultplaceholder & \resultplaceholder \\
\tablelang{World--action base (real-robot base)}{世界—动作基座（真机基座）} & \resultplaceholder & \resultplaceholder & \resultplaceholder & \resultplaceholder & \resultplaceholder & \resultplaceholder & \resultplaceholder \\
\bottomrule
\end{tabular}
\end{table*}

We checked the labels on forty clips, fourteen for development and twenty-six held
out, with five anchors each. The clips cover seven strata: stationary targets,
smooth translation, rotation, contact or slip, occlusion and recovery, weak
appearance, and symmetric or articulated parts. The held-out clips were not used to
tune the estimators. In a manual check of \fillin{200} random episodes, the accepted
box lies on the intended target in \fillin{XX}\% of them. Prefix truncation and
future replacement confirm that emitted inputs, validity, identity, and reference
geometry do not change with future frames. Independent manual landmarks and
multi-view geometry provide references where they are sufficiently constrained,
and we mark physical accuracy as unverified elsewhere. On the held-out clips, the
median translation and rotation errors are \fillin{X}\,mm and \fillin{X}$^\circ$ for
the causal path and \fillin{X}\,mm and \fillin{X}$^\circ$ for the offline path.
Identity failures occur on \fillin{X}\% of the anchors, \fillin{X}\% of the accepted
windows are false accepts, and the labels cover \fillin{XX}\% of the windows. The
check took \fillin{X} hours of annotation, and lifting all of DROID took
\fillin{X} GPU-hours. On surfaces that are static in the image, the selected
PointWorld tracks move by \fillin{0.9} and \fillin{5.8}\,mm (median and 95th
percentile) over the horizons, against \fillin{18} and \fillin{72}\,mm for flow
lifted with monocular depth, whose scale drifts over time. On the moving subset,
where neither source has a reference, we report the agreement of both with a
forward--backward-consistent 2D tracker.

\subsection{Recipe and Scale}
\label{app:scale}
\paragraph{Simulation model.} \ours trains on \bench from the generic visual
initialization with the recipe of Table~\ref{tab:hyper}. We chose its scale and
recipe with a sweep at equal numbers of training samples over two backbone sizes
(the reference and about twice its parameters, by depth at fixed width), two global
batch sizes ($B_0$ and four times $B_0$), and two learning rates per batch size
($\{0.5,1\}\eta_0$ at $B_0$ and $\{1,2\}\eta_0$ at the larger batch), followed by
\fillin{four} local refinements around the best configuration of each size, all
selected on development scenes (Fig.~\ref{fig:scaling}a).

\paragraph{DROID base.} The base uses the same architecture and initialization and
only real data. We trained two pilot sizes, about two and four times the simulation
backbone, for one pass on a stratified quarter of the training windows. We selected
the size by its held-out forecast and action errors, under the deployment budget on
one \cfgGPU{} (p95 time-to-action of at most 200\,ms and no call above 250\,ms). The
selected size (\fillin{XX}M backbone parameters) is pretrained once on all training
windows for \fillin{four} passes at a global batch size of \fillin{256}, with AdamW,
2.5\% warm-up, and cosine decay to 0.1 of the peak learning rate. The checkpoint
with the lowest development loss is the base (Fig.~\ref{fig:scaling}b). We train the
action-only, label-source, and masked-label controls at the smaller pilot size and
budget and compare them with a world--action base of that size
(Table~\ref{tab:supp_pretraining}).

\begin{figure}[t]
\centering
\figplaceholder[1.55in]{\linewidth}{Fig.~S3 --- recipe and scale}{%
\textbf{(a) Simulation sweep.} y: development SR; x groups: the two backbone sizes; within a group, four bars for
$(B_0,0.5\eta_0)$, $(B_0,\eta_0)$, $(4B_0,\eta_0)$, $(4B_0,2\eta_0)$ at equal sample exposure; refinements as small diamonds;
the selected configuration starred.
\textbf{(b) DROID base.} x: training samples (log); y: held-out target translation error at the last horizon (left axis, mm)
and flow endpoint error (right axis, mm); thin lines for the two pilot sizes, thick for the main run; marker at the selected
checkpoint; dotted horizontal line: constant-velocity extrapolation of the same history.}
\caption{Recipe and scale of the simulation model (a) and the DROID base (b).}
\label{fig:scaling}
\end{figure}

\subsection{Independent Post-Training and Controls}
\label{app:posttraining}
The real-robot policy is initialized from the selected base and fine-tuned on the
conveyor demonstrations only. The base was pretrained on detector poses, whereas
the conveyor policy receives poses from the model-based tracker; fine-tuning adapts
the history encoder to the tracker, and we did not measure the cost of this change
of pose source separately. Simulation weights do not initialize
any real-robot policy, and the base does not see simulation data. Each policy
creates its own optimizer, schedule, training-only normalization, and readouts, and
we re-initialize incompatible action or state projections with the same seed in
every control. The reference without the base is fine-tuned from the generic visual
initialization on the same demonstrations, with the same search over \fillin{three}
learning rates and two training lengths on development trials and the same
evaluation. The controlled simulation comparison post-trains the base on \bench with
the downstream recipe and the same data fractions as the generic initialization.
The action-only control pretrains R0 on the same windows, so that its world queries
receive no objective. The label-source control keeps everything except the flow
labels, and the masked-label control removes world labels after the first contact
(Table~\ref{tab:supp_pretraining}).
